% !TEX root = ../../supporting_information.tex

\section{PREDICT Scenario Design and Convergence Assessment}
\label{sec:si-predict-design}

\subsection{Geologic scenario matrix}

\todo{Describe the full factorial construction of the 162 geologic scenarios, including sand proportion, uniform or nonuniform thickness structure, faulting depth, sand clay-volume fraction, and clay clay-volume fraction. State how adjacent units of the same material are collapsed before PREDICT simulation and document the cell-union $P(\mathrm{smear})$ material-placement option.}

\refstepcounter{table}
\label{tab:si-predict-ensemble}
\noindent\textbf{Table \thetable.} Summary of the PREDICT ensemble used by the hierarchical workflow. Replace provisional entries and add the final scenario-level parameter values before submission.
\begin{center}
  \begin{tabular}{@{}ll@{}}
    \toprule
    Item & Current design \\
    \midrule
    Geologic scenarios & 162 \\
    Throw windows per geology & 6 \\
    PREDICT realizations per window & 2000 \\
    Permeability components retained jointly & $k_{xx}$, $k_{yy}$, $k_{zz}$ \\
    Material-placement option & Cell-union $P(\mathrm{smear})$ \\
    Layer representation & Adjacent same-material units collapsed \\
    \bottomrule
  \end{tabular}
\end{center}

\subsection{Convergence assessment}

\todo{Document the three independent 20,000-realization reference ensembles, the reference-distance floor, the repeated smaller ensembles, and the MAE, Hellinger, and Wasserstein diagnostics. Report the quantitative evidence supporting 2000 realizations per geology-window library.}

% Example supporting figure block:
% \begin{figure}[htbp]
%   \centering
%   \includegraphics[width=\textwidth]{predict_convergence_summary.pdf}
%   \caption{Convergence of the empirical permeability distributions toward
%   the reference-distance floor.}
%   \label{fig:si-predict-convergence}
% \end{figure}
